% GENERATED FILE. Edit canonical lessons, then run npm run generate:books.
% canonical-source-hash: fnv1a64:8d21c09bac9bb2d1
% canonical-lessons: UR-C336-sabit-karna, UR-C336-pukhta, UR-C336-ilzam, UR-C336-dastavez, UR-C336-khara

\chapter{Proof}
\label{ch:ur-proof}
\hlchaptermodality{336}

\begin{chapteropening}
\textbf{By the end of this chapter:} \emph{I can ask for proof and give it.}

Complete the last lesson of chapter 336: I can ask for proof and give it.
\end{chapteropening}

\section[\texorpdfstring{\ur{ثابت} \ur{کرنا} (sābit karnā)}{ثابت کرنا (sābit karnā)}]{\ur{ثابت} \ur{کرنا} (sābit karnā) — to prove}

\label{lesson:UR-C336-sabit-karna}



\begin{quote}
Before the new one: say the Urdu for in agreement, then the Urdu for a side in a dispute.
\end{quote}

\subsection*{You'll want to know: \ur{ثابت} \ur{کرنا}}
\textbf{\ur{ثابت} \ur{کرنا}} — \emph{sābit karnā} — \textquotedblleft{}to prove\textquotedblright{}.

\begin{quote}
\ur{آپ} \ur{اپنی} \ur{بات} \ur{کیسے} \ur{ثابت} \ur{کریں} \ur{گے؟}
\end{quote}

\begin{quote}
\emph{āp apnī bāt kaise sābit kareṅ ge?}
\end{quote}

\begin{quote}
How will you prove what you say?
\end{quote}

Past, with \ur{نے} (he or she …): \ur{اس} \ur{نے} \ur{اپنی} \ur{بات} \ur{ثابت} \ur{کی}, \emph{us ne apnī bāt sābit kī}.

\begin{cousinweb}
\textbf{\ur{ثابت} \ur{کرنا}} (\emph{sābit karnā}) is built from \emph{sābit}, usually traced, through Persian, to Arabic \emph{thābit}, \textquotedblleft{}firm, established\textquotedblright{}, and \textbf{\ur{کرنا}}, \textquotedblleft{}to do\textquotedblright{}: to make a claim stand firm.
\end{cousinweb}

\begin{grammarlens}[title={What you've built}]
One more word for giving reasons and opinions.
\end{grammarlens}

\subsection*{Your turn}
\begin{itemize}
\raggedright
  \item \emph{Say it:} \emph{sābit karnā}
  \item \emph{Say it:} \emph{sābit karnā}, once more
  \item \emph{Say it:} \emph{āp apnī bāt kaise sābit kareṅ ge?} — How will you prove what you say?
  \item \emph{Recall:} say the Urdu for in agreement, then the Urdu for a side in a dispute, then say \emph{sābit karnā} again
\end{itemize}

\begin{tcolorbox}[breakable,colback=teal!4,colframe=teal!35!black,title={Before you move on}]
What does \ur{ثابت} \ur{کرنا} mean? (\textquotedblleft{}To prove\textquotedblright{}.) Say it once more.
\end{tcolorbox}
\section[\texorpdfstring{\ur{پختہ} (pukhtā)}{پختہ (pukhtā)}]{\ur{پختہ} (pukhtā) — firm; conclusive}

\label{lesson:UR-C336-pukhta}



\begin{quote}
Before the new one: say the Urdu for a side in a dispute, then the Urdu for to prove.
\end{quote}

\subsection*{You'll want to know: \ur{پختہ}}
\textbf{\ur{پختہ}} — \emph{pukhtā} — \textquotedblleft{}firm; conclusive\textquotedblright{}.

\begin{quote}
\ur{اس} \ur{کا} \ur{ارادہ} \ur{پختہ} \ur{ہے۔}
\end{quote}

\begin{quote}
\emph{us kā irādā pukhtā hai.}
\end{quote}

\begin{quote}
Her mind is firmly made up.
\end{quote}

\begin{cousinweb}
\textbf{\ur{پختہ}} (\emph{pukhtā}) is usually traced to Persian \emph{pukhta}, \textquotedblleft{}cooked, ripe\textquotedblright{}: what is fully done is firm.
\end{cousinweb}

\begin{grammarlens}[title={What you've built}]
One more word for giving reasons and opinions.
\end{grammarlens}

\subsection*{Your turn}
\begin{itemize}
\raggedright
  \item \emph{Say it:} \emph{pukhtā}
  \item \emph{Say it:} \emph{pukhtā}, once more
  \item \emph{Say it:} \emph{us kā irādā pukhtā hai.} — Her mind is firmly made up.
  \item \emph{Recall:} say the Urdu for a side in a dispute, then the Urdu for to prove, then say \emph{pukhtā} again
\end{itemize}

\begin{tcolorbox}[breakable,colback=teal!4,colframe=teal!35!black,title={Before you move on}]
What does \ur{پختہ} mean? (\textquotedblleft{}Firm; conclusive\textquotedblright{}.) Say it once more.
\end{tcolorbox}
\section[\texorpdfstring{\ur{الزام} (ilzām)}{الزام (ilzām)}]{\ur{الزام} (ilzām) — an accusation; a charge}

\label{lesson:UR-C336-ilzam}



\begin{quote}
Before the new one: say the Urdu for to prove, then the Urdu for firm.
\end{quote}

\subsection*{You'll want to know: \ur{الزام}}
\textbf{\ur{الزام}} — \emph{ilzām} — \textquotedblleft{}an accusation; a charge\textquotedblright{}.

\begin{quote}
\ur{یہ} \ur{الزام} \ur{بالکل} \ur{غلط} \ur{ہے۔}
\end{quote}

\begin{quote}
\emph{yih ilzām bilkul ghalat hai.}
\end{quote}

\begin{quote}
This accusation is completely wrong.
\end{quote}

Masculine.

\begin{cousinweb}
\textbf{\ur{الزام}} (\emph{ilzām}) is usually traced, through Persian, to Arabic \emph{ilzām}, \textquotedblleft{}binding someone to something\textquotedblright{}; the sense \textquotedblleft{}a charge laid on someone\textquotedblright{} is the one Urdu uses.
\end{cousinweb}

\begin{grammarlens}[title={What you've built}]
One more word for giving reasons and opinions.
\end{grammarlens}

\subsection*{Your turn}
\begin{itemize}
\raggedright
  \item \emph{Say it:} \emph{ilzām}
  \item \emph{Say it:} \emph{ilzām}, once more
  \item \emph{Say it:} \emph{yih ilzām bilkul ghalat hai.} — This accusation is completely wrong.
  \item \emph{Recall:} say the Urdu for to prove, then the Urdu for firm, then say \emph{ilzām} again
\end{itemize}

\begin{tcolorbox}[breakable,colback=teal!4,colframe=teal!35!black,title={Before you move on}]
What does \ur{الزام} mean? (\textquotedblleft{}An accusation; a charge\textquotedblright{}.) Say it once more.
\end{tcolorbox}
\section[\texorpdfstring{\ur{دستاویز} (dastāvez)}{دستاویز (dastāvez)}]{\ur{دستاویز} (dastāvez) — a document}

\label{lesson:UR-C336-dastavez}



\begin{quote}
Before the new one: say the Urdu for firm, then the Urdu for an accusation.
\end{quote}

\subsection*{You'll want to know: \ur{دستاویز}}
\textbf{\ur{دستاویز}} — \emph{dastāvez} — \textquotedblleft{}a document\textquotedblright{}.

\begin{quote}
\ur{یہ} \ur{دستاویز} \ur{ہمارے} \ur{دفتر} \ur{میں} \ur{ہے۔}
\end{quote}

\begin{quote}
\emph{yih dastāvez hamāre daftar meṅ hai.}
\end{quote}

\begin{quote}
This document is in our office.
\end{quote}

Feminine.

\begin{cousinweb}
\textbf{\ur{دستاویز}} (\emph{dastāvez}) is usually traced to Persian \emph{dastāwēz}, \textquotedblleft{}a document, a deed\textquotedblright{}.
\end{cousinweb}

\begin{grammarlens}[title={What you've built}]
One more word for giving reasons and opinions.
\end{grammarlens}

\subsection*{Your turn}
\begin{itemize}
\raggedright
  \item \emph{Say it:} \emph{dastāvez}
  \item \emph{Say it:} \emph{dastāvez}, once more
  \item \emph{Say it:} \emph{yih dastāvez hamāre daftar meṅ hai.} — This document is in our office.
  \item \emph{Recall:} say the Urdu for firm, then the Urdu for an accusation, then say \emph{dastāvez} again
\end{itemize}

\begin{tcolorbox}[breakable,colback=teal!4,colframe=teal!35!black,title={Before you move on}]
What does \ur{دستاویز} mean? (\textquotedblleft{}A document\textquotedblright{}.) Say it once more.
\end{tcolorbox}
\section[\texorpdfstring{\ur{کھرا} (kharā)}{کھرا (kharā)}]{\ur{کھرا} (kharā) — genuine; straight (honest)}

\label{lesson:UR-C336-khara}



\begin{quote}
Before the new one: say the Urdu for an accusation, then the Urdu for a document.
\end{quote}

\subsection*{You'll want to know: \ur{کھرا}}
\textbf{\ur{کھرا}} — \emph{kharā} — \textquotedblleft{}genuine; straight (honest)\textquotedblright{}.

\begin{quote}
\ur{وہ} \ur{کھرا} \ur{آدمی} \ur{ہے،} \ur{جھوٹ} \ur{نہیں} \ur{بولتا۔}
\end{quote}

\begin{quote}
\emph{voh kharā ādmī hai, jhūṭ nahīṅ boltā.}
\end{quote}

\begin{quote}
He is a straight man; he does not lie.
\end{quote}

\begin{cousinweb}
\textbf{\ur{کھرا}} (\emph{kharā}) is inherited: CDIAL lists the Hindi cousin \emph{kharā}, \textquotedblleft{}sharp, pure, good\textquotedblright{}, under Sanskrit \emph{khára}, \textquotedblleft{}hard, sharp\textquotedblright{} (CDIAL 3819).
\end{cousinweb}

\begin{grammarlens}[title={What you've built}]
One more word for giving reasons and opinions.
\end{grammarlens}

\subsection*{Your turn}
\begin{itemize}
\raggedright
  \item \emph{Say it:} \emph{kharā}
  \item \emph{Say it:} \emph{kharā}, once more
  \item \emph{Say it:} \emph{voh kharā ādmī hai, jhūṭ nahīṅ boltā.} — He is a straight man; he does not lie.
  \item \emph{Recall:} say the Urdu for an accusation, then the Urdu for a document, then say \emph{kharā} again
\end{itemize}

\begin{tcolorbox}[breakable,colback=teal!4,colframe=teal!35!black,title={Before you move on}]
What does \ur{کھرا} mean? (\textquotedblleft{}Genuine; straight (honest)\textquotedblright{}.) Say it once more.
\end{tcolorbox}
